\documentclass[10pt]{article}
\usepackage[T1]{fontenc}
\usepackage[utf8]{inputenc}
\usepackage[french]{babel}
\usepackage{amsmath,amssymb,mathrsfs,mathpazo}
\usepackage{geometry}
\geometry{a4paper,margin=2cm}
\usepackage{graphicx,array,tabularx,tabularray,colortbl,multirow,multicol,booktabs,enumitem,pifont,ulem,fancybox,diagbox,cancel}
\usepackage{pstricks,pstricks-add,pst-plot,pst-tree,pst-node,pst-text,pst-func,pst-math,pst-bspline,pst-3dplot,pst-eucl,pst-solides3d}
\let\pscalculate\undefined
\usepackage{pst-bezier}
\usepackage{tikz,pgfplots,tkz-tab}
\usetikzlibrary{arrows,arrows.meta,patterns,calc,positioning,trees,decorations.markings,shapes.geometric,shapes.misc,shapes.arrows,intersections,decorations.pathmorphing,decorations.pathreplacing,backgrounds,shadows,babel}
\pgfplotsset{compat=1.18}
\usepgfplotslibrary{fillbetween}
\usepackage[np]{numprint}
\usepackage{listings,makecell,pas-tableur,twoopt}
\let\barre\undefined
\usepackage[pstricks]{tablvar}
\definecolor{darkgreen}{RGB}{0,100,0}
\providecommand{\red}{\color{red}}
\providecommand{\blue}{\color{blue}}
\providecommand{\green}{\color{green}}
\providecommand{\black}{\color{black}}
\providecommand{\white}{\color{white}}
\providecommand{\gray}{\color{gray}}
\providecommand{\cyan}{\color{cyan}}
\providecommand{\magenta}{\color{magenta}}
\providecommand{\yellow}{\color{yellow}}
\providecommand{\violet}{\color{violet}}
\providecommand{\orange}{\color{orange}}
\providecommand{\euro}{\texteuro}
\providecommand{\EUR}{\texteuro}
\providecommand{\ieme}{\textsuperscript{e}}
\providecommand{\degres}{\ensuremath{^\circ}}
\providecommand{\degree}{\ensuremath{^\circ}}
\providecommand{\vv}[1]{\overrightarrow{#1}}
\providecommand{\Coord}[1]{#1}
\providecommand{\decosix}{}
\providecommand{\starredbullet}{\ensuremath{\bullet}}
\pagestyle{empty}
\setlength{\parindent}{0pt}
\usepackage{adjustbox}\DeclareUnicodeCharacter{2236}{:}\DeclareUnicodeCharacter{2190}{\ensuremath{\leftarrow}}
\providecommand{\R}{\mathbb{R}}
\providecommand{\N}{\mathbb{N}}
\providecommand{\Z}{\mathbb{Z}}
\providecommand{\Q}{\mathbb{Q}}
\providecommand{\D}{\mathbb{D}}
\providecommand{\C}{\mathbb{C}}
\providecommand{\e}{{\mathrm{e}}}
\providecommand{\vect}[1]{\overrightarrow{#1}}
\providecommand{\vectt}[1]{\overrightarrow{\mathrm{#1}}}
\providecommand{\vvt}[1]{\overrightarrow{\mathrm{#1}}}
\providecommand{\np}[2][]{{#2~#1}}
\providecommand{\vv}[1]{{\overrightarrow{#1}}}
\renewcommand{\d}{{\mathrm{d}}}
\providecommand{\degres}{{^\circ}}
\providecommand{\degre}{{^\circ}}
\providecommand{\blue}{\color{blue}}
\providecommand{\red}{\color{red}}
\providecommand{\enskip}{\quad}
\providecommand{\llbracket}{\lbrack\!\lbrack}
\providecommand{\rrbracket}{\rbrack\!\rbrack}
\providecommand{\euro}{{\text{€}}}
\providecommand{\no}{{\text{n°}}}
\providecommand{\interval}[2]{{[#1;#2]}}
\providecommand{\intervFO}[2]{{[#1;#2[}}
\providecommand{\dotfill}{\ldots}

\providecommand{\ds}{\displaystyle}
\providecommand{\textding}[1]{#1}
\providecommand{\Oij}{\(\mathrm{O};\vec{\imath},\vec{\jmath}\)}
\providecommand{\Oijk}{\(\mathrm{O};\vec{\imath},\vec{\jmath},\vec{k}\)}
\providecommand{\Ouv}{\(\mathrm{O};\vec{u},\vec{v}\)}
\providecommand{\intervFF}[2]{\ensuremath{[#1;#2]}}
\renewcommand{\euro}{\ifmmode\text{\texteuro}\else\texteuro\fi}
\DeclareUnicodeCharacter{00B0}{\ifmmode^{\circ}\else\textdegree\fi}
\DeclareUnicodeCharacter{20AC}{\ifmmode\text{\texteuro}\else\texteuro\fi}
% Énoncés : Éducation nationale. Transcription et corrigés : APMEP, Denis Vergès et contributeurs. Découpage et classement : CoopMaths, Sébastien Lozano, Éric Elter et contributeurs. Adaptation de la présentation : Maths Explorer.
% CC BY-SA : https://coopmaths.fr/www/about/
% https://coopmaths.fr/alea/static/eam/2026/tex/eam_specialite_2026_06_etrangers_1_cor.tex
\begin{document}
\noindent\textit{APMEP / CoopMaths. CC BY-SA. Présentation adaptée par Maths Explorer.}\par\medskip
Un maire souhaite végétaliser sa ville. Pour cela, il décide de planter des mûriers platanes dans les différents parcs.
			
			Ces arbres sont réputés pour leurs qualités d'ombrage et de résistance à la sécheresse.
			
			\subsection*{Partie A}
			
			Au moment de la plantation, un mûrier platane mesure $1$ mètre.
			
			On suppose que chaque année la hauteur de l'arbre augmente de $40$ cm, c'est-à-dire de $0{,}40$ m.
			
			Pour tout entier naturel $n$, on note $u_n$ la hauteur de l'arbre, en mètres, $n$ années après sa plantation. Ainsi $u_0=1$.
			
			\begin{enumerate}
				\item \begin{enumerate}
					\item Une année après la plantation, la hauteur de l'arbre a augmenté de $0{,}40$ m.
					
					On a donc :
					\begin{align*}
						u_1&=u_0+0{,}40\\
						&=1+0{,}40\\
						&=1{,}40.
					\end{align*}
					
					Ainsi,
					\[
					u_1=1{,}40.
					\]
					
					Un an après la plantation, l'arbre mesure donc $1{,}40$ m.
					
					\item Deux années après la plantation, l'arbre a grandi deux fois de $0{,}40$ m.
					
					On peut calculer :
					\begin{align*}
						u_2&=u_1+0{,}40\\
						&=1{,}40+0{,}40\\
						&=1{,}80.
					\end{align*}
					
					Ainsi, deux années après sa plantation, l'arbre mesure $1{,}80$ m.
				\end{enumerate}
				
				\item Chaque année, la hauteur de l'arbre augmente d'une même quantité, égale à $0{,}40$ m.
				
				Ainsi, pour tout entier naturel $n$, on a :
				\[
				u_{n+1}=u_n+0{,}40.
				\]
				
				La suite $(u_n)$ est donc une suite arithmétique de raison $0{,}40$ et de premier terme $u_0=1$.
				
				\item Comme $(u_n)$ est une suite arithmétique de premier terme $u_0=1$ et de raison $0{,}40$, on a, pour tout entier naturel $n$ :
				\[
				u_n=u_0+n\times 0{,}40.
				\]
				
				Donc, pour tout entier naturel $n$ :
				\[
				u_n=1+0{,}40n.
				\]
				
				\item On cherche au bout de combien d'années le mûrier atteindra $9$ mètres de haut.
				
				Il faut donc résoudre l'inéquation :
				\[
				u_n\geq 9.
				\]
				
				Or, pour tout entier naturel $n$, $u_n=1+0{,}40n$. Ainsi :
				\begin{align*}
					&\phantom{\iff}u_n\geq 9\\
					&\iff 1+0{,}40n\geq 9\\
					&\iff 0{,}40n\geq 8\\
					&\iff n\geq \frac{8}{0{,}40}\\
					&\iff n\geq 20.
				\end{align*}
				
				Le mûrier atteindra donc $9$ mètres de haut au bout de $20$ ans.
			\end{enumerate}
			
			\subsection*{Partie B}
			
			Au moment de la plantation, l'arbre possède un tronc et deux branches.
			
			Un an après la plantation, on observe $4$ nouvelles branches.
			
			Deux ans après la plantation, on observe $8$ nouvelles branches.
			
			Chaque année, le nombre de nouvelles branches double comme représenté sur le schéma ci-dessous :
			
			\begin{center}
				\begin{tikzpicture}[
					scale=0.82,
					point/.style={circle, draw=black, fill=gray!65, inner sep=1.8pt},
					branche/.style={line width=1.3pt, gray!75!black},
					sol/.style={line width=1.2pt, teal!70!black}
					]
					
					
					\draw[sol] (-1.7,0) -- (12.7,0);
					
					
					\coordinate (T0) at (1,0);
					\coordinate (N0) at (1,1.45);
					\coordinate (G0) at (-0.15,2.05);
					\coordinate (D0) at (2.15,2.05);
					
					\draw[branche] (T0) -- (N0);
					\draw[branche] (N0) -- (G0);
					\draw[branche] (N0) -- (D0);
					
					\node[point] at (T0) {};
					\node[point] at (N0) {};
					\node[point] at (G0) {};
					\node[point] at (D0) {};
					
					\node[left] at (-0.35,1.5) {branche};
					\node[right] at (1.7,1.5) {branche};
					\node[left] at (0.65,0.45) {tronc};
					
					\node[align=center] at (1,-1) {au moment\\de la plantation\\(année 0)};
					
					
					\coordinate (T1) at (5.2,0);
					\coordinate (N1) at (5.2,1.45);
					\coordinate (G1) at (4.05,2.05);
					\coordinate (D1) at (6.35,2.05);
					\coordinate (GG1) at (3.5,3.1);
					\coordinate (GD1) at (4.6,3.1);
					\coordinate (DG1) at (5.8,3.1);
					\coordinate (DD1) at (6.9,3.1);
					
					\draw[branche] (T1) -- (N1);
					\draw[branche] (N1) -- (G1);
					\draw[branche] (N1) -- (D1);
					\draw[branche] (G1) -- (GG1);
					\draw[branche] (G1) -- (GD1);
					\draw[branche] (D1) -- (DG1);
					\draw[branche] (D1) -- (DD1);
					
					\foreach \P in {T1,N1,G1,D1,GG1,GD1,DG1,DD1}
					{
						\node[point] at (\P) {};
					}
					
					\node[align=center] at (5.2,-1) {un an après\\la plantation\\(année 1)};
					
					
					\coordinate (T2) at (10.1,0);
					\coordinate (N2) at (10.1,1.6);
					\coordinate (G2) at (8.85,2.3);
					\coordinate (D2) at (11.35,2.3);
					\coordinate (GG2) at (8.2,3.25);
					\coordinate (GD2) at (9.15,3.25);
					\coordinate (DG2) at (10.85,3.25);
					\coordinate (DD2) at (11.8,3.25);
					\coordinate (A2) at (7.8,4.2);
					\coordinate (B2) at (8.55,4.2);
					\coordinate (C2) at (8.9,4.2);
					\coordinate (D2top) at (9.65,4.2);
					\coordinate (E2) at (10.45,4.2);
					\coordinate (F2) at (11.2,4.2);
					\coordinate (G2top) at (11.55,4.2);
					\coordinate (H2) at (12.3,4.2);
					
					\draw[branche] (T2) -- (N2);
					\draw[branche] (N2) -- (G2);
					\draw[branche] (N2) -- (D2);
					\draw[branche] (G2) -- (GG2);
					\draw[branche] (G2) -- (GD2);
					\draw[branche] (D2) -- (DG2);
					\draw[branche] (D2) -- (DD2);
					\draw[branche] (GG2) -- (A2);
					\draw[branche] (GG2) -- (B2);
					\draw[branche] (GD2) -- (C2);
					\draw[branche] (GD2) -- (D2top);
					\draw[branche] (DG2) -- (E2);
					\draw[branche] (DG2) -- (F2);
					\draw[branche] (DD2) -- (G2top);
					\draw[branche] (DD2) -- (H2);
					
					\foreach \P in {T2,N2,G2,D2,GG2,GD2,DG2,DD2,A2,B2,C2,D2top,E2,F2,G2top,H2}
					{
						\node[point] at (\P) {};
					}
					
					\node[align=center] at (10.1,-1) {deux ans après\\la plantation\\(année 2)};
					
				\end{tikzpicture}
			\end{center}
			
			Pour tout entier naturel $n$, on note $v_n$ le nombre de nouvelles branches $n$ années après la plantation. À la plantation, l'arbre possède $2$ branches, ainsi on pose $v_0=2$.
			
			\begin{enumerate}
				\item Chaque année, le nombre de nouvelles branches double.
				
				Ainsi, pour tout entier naturel $n$, on a :
				\[
				v_{n+1}=2v_n.
				\]
				
				La suite $(v_n)$ est donc une suite géométrique de raison $2$ et de premier terme $v_0=2$.
				
				On peut en déduire que, pour tout entier naturel $n$ :
				\[
				v_n=v_0\times 2^n.
				\]
				
				Donc :
				\[
				v_n=2\times 2^n=2^{n+1}.
				\]
				
				\item \begin{enumerate}
					\item Trois ans après la plantation, l'arbre possède toutes les branches déjà présentes à la plantation, ainsi que les nouvelles branches produites les années $1$, $2$ et $3$.
					
					On calcule d'abord le nombre de nouvelles branches produites chaque année :
					\[
					v_0=2,\qquad v_1=4,\qquad v_2=8,\qquad v_3=16.
					\]
					
					Le nombre total de branches trois ans après la plantation est donc :
					\begin{align*}
						v_0+v_1+v_2+v_3
						&=2+4+8+16\\
						&=30.
					\end{align*}
					
					Ainsi, trois ans après sa plantation, l'arbre possède bien $30$ branches au total.
					
					\item On considère le programme Python suivant :
					
					\begin{center}
						\fbox{\begin{minipage}{0.42\linewidth}\ttfamily
								v = 2\par
								total = 2\par
								for i in range(10):\par
								\quad v = 2 * v\par
								\quad total = total + v\par
								print(total)
						\end{minipage}}
					\end{center}
					
					Au début du programme, on a :
					\[
					v=2 \qquad \text{et} \qquad \text{total}=2.
					\]
					
					La variable $v$ représente le nombre de nouvelles branches produites pendant une année donnée.
					
					La variable $\text{total}$ représente le nombre total de branches de l'arbre depuis la plantation jusqu'à l'année considérée.
					
					La boucle \texttt{for i in range(10)} répète $10$ fois les instructions suivantes :
					\begin{itemize}[left=5mm,label={\textbullet~~}]
						\item on double la valeur de $v$ ;
						\item on ajoute cette nouvelle valeur de $v$ au total.
					\end{itemize}
					
					Après $10$ répétitions, le programme a donc ajouté les nombres :
					\[
					v_1,\ v_2,\ v_3,\ \ldots,\ v_{10}
					\]
					
					au nombre initial de branches $v_0=2$.
					
					La valeur affichée est donc :
					\begin{align*}
						v_0+v_1+v_2+\cdots+v_{10}
						&=2+4+8+\cdots+2048\\
						&=4094.
					\end{align*}
					
					Dans le contexte de l'exercice, la valeur $4\ 094$ représente donc le nombre total de branches que possède l'arbre $10$ ans après sa plantation.
				\end{enumerate}
			\end{enumerate}

\clearpage

\end{document}
